\section{الشكل المصفوفي العام للنموذج الديناميكي للذراع }


بما أن الذراع الروبوتية المدروسة تتكون من خمس درجات حرية، فإن عدد متحولات المفاصل يساوي:

\begin{equation}
n = 5
\end{equation}

وبالتالي يمكن كتابة متجه المتحولات المفصلية بالشكل:

\begin{equation}
q =
\begin{bmatrix}
\theta_1 &
\theta_2 &
\theta_3 &
\theta_4 &
\theta_5
\end{bmatrix}^{T}
\end{equation}

كما يمكن كتابة متجه السرعات المفصلية بالشكل:

\begin{equation}
\dot{q} =
\begin{bmatrix}
\dot{\theta}_1 &
\dot{\theta}_2 &
\dot{\theta}_3 &
\dot{\theta}_4 &
\dot{\theta}_5
\end{bmatrix}^{T}
\end{equation}

أما متجه التسارعات المفصلية فيكتب كما يلي:

\begin{equation}
\ddot{q} =
\begin{bmatrix}
\ddot{\theta}_1 &
\ddot{\theta}_2 &
\ddot{\theta}_3 &
\ddot{\theta}_4 &
\ddot{\theta}_5
\end{bmatrix}^{T}
\end{equation}

ويمثل متجه العزوم المطبقة على مفاصل الذراع بالشكل:

\begin{equation}
\tau =
\begin{bmatrix}
\tau_1 \\
\tau_2 \\
\tau_3 \\
\tau_4 \\
\tau_5
\end{bmatrix}
\end{equation}

وبالتعويض عن عدد درجات الحرية $n=5$ في النموذج الديناميكي العام الناتج عن صياغة لاغرانج، يمكن كتابة النموذج المصفوفي للذراع الروبوتية كما يلي:

\begin{equation}
\tau =
\left(
D(q) + I_{act}
\right)\ddot{q}
+
D_{cen}(q,\dot{q})
+
D_{cor}(q,\dot{q})
+
D_g(q)
\end{equation}

حيث تمثل $D(q)$ مصفوفة العطالة، وتمثل $I_{act}$ مصفوفة عزوم القصور الذاتي للمشغلات، بينما تمثل $D_{cen}$ حدود القوى الطاردة المركزية، وتمثل $D_{cor}$ حدود قوى كوريوليس، أما $D_g$ فيمثل متجه تأثير الجاذبية.

\subsection*{مصفوفة العطالة}

تعطى مصفوفة العطالة للذراع الروبوتية ذات خمس درجات حرية بالشكل:

\begin{equation}
D(q)=
\begin{bmatrix}
D_{11} & D_{12} & D_{13} & D_{14} & D_{15} \\
D_{21} & D_{22} & D_{23} & D_{24} & D_{25} \\
D_{31} & D_{32} & D_{33} & D_{34} & D_{35} \\
D_{41} & D_{42} & D_{43} & D_{44} & D_{45} \\
D_{51} & D_{52} & D_{53} & D_{54} & D_{55}
\end{bmatrix}
\end{equation}

وتؤثر هذه المصفوفة في متجه التسارعات المفصلية وفق العلاقة:

\begin{equation}
D(q)\ddot{q} =
\begin{bmatrix}
D_{11} & D_{12} & D_{13} & D_{14} & D_{15} \\
D_{21} & D_{22} & D_{23} & D_{24} & D_{25} \\
D_{31} & D_{32} & D_{33} & D_{34} & D_{35} \\
D_{41} & D_{42} & D_{43} & D_{44} & D_{45} \\
D_{51} & D_{52} & D_{53} & D_{54} & D_{55}
\end{bmatrix}
\begin{bmatrix}
\ddot{\theta}_1 \\
\ddot{\theta}_2 \\
\ddot{\theta}_3 \\
\ddot{\theta}_4 \\
\ddot{\theta}_5
\end{bmatrix}
\end{equation}

\subsection*{مصفوفة عزوم القصور الذاتي للمشغلات}

تؤخذ عزوم القصور الذاتي للمشغلات بعين الاعتبار من خلال مصفوفة قطرية، وذلك لأن كل مشغل يؤثر بشكل مباشر على المفصل المرتبط به. وتكتب هذه المصفوفة بالشكل:

\begin{equation}
I_{act} =
\begin{bmatrix}
I_{1(act)} & 0 & 0 & 0 & 0 \\
0 & I_{2(act)} & 0 & 0 & 0 \\
0 & 0 & I_{3(act)} & 0 & 0 \\
0 & 0 & 0 & I_{4(act)} & 0 \\
0 & 0 & 0 & 0 & I_{5(act)}
\end{bmatrix}
\end{equation}

وبالتالي يكون تأثير قصور المشغلات على التسارعات المفصلية كما يلي:

\begin{equation}
I_{act}\ddot{q} =
\begin{bmatrix}
I_{1(act)} & 0 & 0 & 0 & 0 \\
0 & I_{2(act)} & 0 & 0 & 0 \\
0 & 0 & I_{3(act)} & 0 & 0 \\
0 & 0 & 0 & I_{4(act)} & 0 \\
0 & 0 & 0 & 0 & I_{5(act)}
\end{bmatrix}
\begin{bmatrix}
\ddot{\theta}_1 \\
\ddot{\theta}_2 \\
\ddot{\theta}_3 \\
\ddot{\theta}_4 \\
\ddot{\theta}_5
\end{bmatrix}
\end{equation}

\subsection*{حدود القوى الطاردة المركزية}

تنتج حدود القوى الطاردة المركزية عن مربعات السرعات الزاوية للمفاصل، ويمكن كتابتها بالشكل المصفوفي التالي:

\begin{equation}
D_{cen}(q,\dot{q}) =
\begin{bmatrix}
D_{111} & D_{122} & D_{133} & D_{144} & D_{155} \\
D_{211} & D_{222} & D_{233} & D_{244} & D_{255} \\
D_{311} & D_{322} & D_{333} & D_{344} & D_{355} \\
D_{411} & D_{422} & D_{433} & D_{444} & D_{455} \\
D_{511} & D_{522} & D_{533} & D_{544} & D_{555}
\end{bmatrix}
\begin{bmatrix}
\dot{\theta}_1^2 \\
\dot{\theta}_2^2 \\
\dot{\theta}_3^2 \\
\dot{\theta}_4^2 \\
\dot{\theta}_5^2
\end{bmatrix}
\end{equation}

\subsection*{حدود قوى كوريوليس}

تنتج حدود كوريوليس والقوى النابذة المركزية عن جداءات السرعات الزاوية للمفاصل، وتظهر هذه الحدود في النموذج الديناميكي كما يلي:
\begin{equation*}
\resizebox{\textwidth}{!}{$
D_{cor}(q,\dot{q}) =
\begin{bmatrix}
D_{112} & D_{113} & D_{114} & D_{115} & D_{123} & D_{124} & D_{125} & D_{134} & D_{135} & D_{145} \\
D_{212} & D_{213} & D_{214} & D_{215} & D_{223} & D_{224} & D_{225} & D_{234} & D_{235} & D_{245} \\
D_{312} & D_{313} & D_{314} & D_{315} & D_{323} & D_{324} & D_{325} & D_{334} & D_{335} & D_{345} \\
D_{412} & D_{413} & D_{414} & D_{415} & D_{423} & D_{424} & D_{425} & D_{434} & D_{435} & D_{445} \\
D_{512} & D_{513} & D_{514} & D_{515} & D_{523} & D_{524} & D_{525} & D_{534} & D_{535} & D_{545}
\end{bmatrix}
\begin{bmatrix}
\dot{\theta}_1\dot{\theta}_2 \\
\dot{\theta}_1\dot{\theta}_3 \\
\dot{\theta}_1\dot{\theta}_4 \\
\dot{\theta}_1\dot{\theta}_5 \\
\dot{\theta}_2\dot{\theta}_3 \\
\dot{\theta}_2\dot{\theta}_4 \\
\dot{\theta}_2\dot{\theta}_5 \\
\dot{\theta}_3\dot{\theta}_4 \\
\dot{\theta}_3\dot{\theta}_5 \\
\dot{\theta}_4\dot{\theta}_5
\end{bmatrix}
$}
\end{equation*}

\subsection*{متجه الجاذبية}

يمثل متجه الجاذبية العزوم الناتجة عن تأثير وزن الوصلات على مفاصل الذراع الروبوتية، ويكتب بالشكل:

\begin{equation}
D_g(q)=
\begin{bmatrix}
D_1 \\
D_2 \\
D_3 \\
D_4 \\
D_5
\end{bmatrix}
\end{equation}

وبذلك يمكن كتابة النموذج الديناميكي الكامل للذراع الروبوتية ذات خمس درجات حرية على الصورة النهائية التالية:

\begin{equation}
\begin{bmatrix}
\tau_1 \\
\tau_2 \\
\tau_3 \\
\tau_4 \\
\tau_5
\end{bmatrix}
=
\left(
D(q)+I_{act}
\right)
\begin{bmatrix}
\ddot{\theta}_1 \\
\ddot{\theta}_2 \\
\ddot{\theta}_3 \\
\ddot{\theta}_4 \\
\ddot{\theta}_5
\end{bmatrix}
+
D_{cen}(q,\dot{q})
+
D_{cor}(q,\dot{q})
+
\begin{bmatrix}
D_1 \\
D_2 \\
D_3 \\
D_4 \\
D_5
\end{bmatrix}
\end{equation}

كما يمكن التعبير عن النموذج الديناميكي بالصورة القياسية الأكثر استخدامًا في نمذجة الأذرع الروبوتية كما يلي:

\begin{equation}
\tau =
M(q)\ddot{q}
+
C(q,\dot{q})\dot{q}
+
G(q)
\end{equation}

حيث إن:

\begin{itemize}
    \item $M(q)$ تمثل مصفوفة العطالة الكلية للذراع مع تأثير المشغلات.
    \item $C(q,\dot{q})\dot{q}$ تمثل حدود قوى كوريوليس والقوى الطاردة المركزية.
    \item $G(q)$ تمثل متجه الجاذبية.
\end{itemize}

